\section{Sylow-equivariant Frobenius comparison}\label{sec:periodicity}

We prove the comparison property used in the extension argument.  It is
important to retain the specified inner factors of the action: a comparison
of the underlying block algebras alone would not control the crossed
multiplications in Section~\ref{sec:extensions}.

\begin{proposition}[Sylow comparison]\label{prop:sylow-comparison}
Fix the coefficient prime $p$ and the defect-order bound $M$.
For the comparison data of Definition~\ref{def:ext:comparison-data}, let
$S\leq Y_0$ be a $p$-subgroup contained in the pure $X$-actions.  There is a
positive integer $m$, bounded in terms of $p$ and $M$, and a Morita
bimodule
\[
 \mathcal M:\quad
 \cO\bar N\sigma^m(\bar b)\text{-}\cO\bar N\bar b
\]
such that its reduction $M_0=k\otimes_{\cO}\mathcal M$ has invertible
$k$-linear maps $J_x$, for $x\in S$, with $J_1=\id$, satisfying
\begin{align}
 J_x(avb)&=\alpha'_x(a)J_x(v)\alpha_x(b),
       \label{per:covariance}\\
 J_xJ_y(v)&=a'_{xy}J_{xy}(v)a_{xy}^{-1}.
       \label{per:factor-law}
\end{align}
Here $\alpha_x$ is the action of the chosen representative $h_x$,
$\alpha'_x=\sigma^m(\alpha_x)$,
$a_{xy}=n_{xy}\bar b$, and
$a'_{xy}=n_{xy}\sigma^m(\bar b)=\sigma^m(a_{xy})$.
The assertion includes $S=1$.  The bound does not depend on the orders of
the universal-cover kernels or of the added central tori.
\end{proposition}

We use the block idempotent and the corresponding block algebra
interchangeably when specifying bimodule sides.  Extend Witt Frobenius to
an algebraic closure of $\operatorname{Frac}(\cO)$ by requiring it to fix
all $p$-power roots of unity.  Such an extension exists because the
cyclotomic extensions generated by these roots are totally ramified,
whereas $\cO$ is unramified; see \cite[Lemma~2.2]{FarrellKessar}.
On roots of unity of order prime to $p$ it is the $p$th power map.
Throughout this section a bound is allowed to depend on $p,M$.

\subsection{Replacing field automorphisms by algebraic permutations}

We use the component models of Lemma~\ref{lem:ext:partial-extension}.
Lift the block to the central product cover, and on a Lie component adjoin
the direct central $p$-factor $C_i$ so that the finite group is the full
fixed-point group $\mathbf J_i^{F_i}$.  These lifted blocks can have
unbounded defect.  We will descend the equivalences before using any
bounded-defect finiteness criterion.

\begin{lemma}[Algebraic realization]\label{per:algebraic-action}
On each $S$-orbit of cross-characteristic Lie components there is a
connected reductive group $\mathbf J$ with simply connected derived
group and a Steinberg endomorphism $F$, together with an identification
of $J=\mathbf J^F$ with the product of the full component groups, having
the following properties.
The operations generated by the $h_x$ and the inner automorphisms of
$J$ form a subgroup
\[
 I\ \leq\ \mathbf J_{\mathrm{ad}}^F\rtimes Y',
\]
where $Y'$ is a finite $p$-group of algebraic automorphisms commuting with
$F$.  Its order and the number of copies of each original absolute
component used in $\mathbf J$ are bounded.  Restriction of these
operations to the universal components is faithful modulo precisely
the standard inner and pinned relations.  In particular the products
of the chosen operations have the prescribed inner factors coming from
$n_{xy}$.
\end{lemma}

\begin{proof}
We recall the relevant part of Steinberg's automorphism theorem
\cite{SteinbergAutomorphisms} for finite simple groups of Lie type,
in the formulation of Broto--M\o ller--Oliver
\cite[Definition~3.3 and Theorem~3.4]{BMO}, omitting the finite exceptional
list already separated in Section~\ref{sec:extensions}.  For an ordinary
diagram twist
\[
 F_0=\gamma\operatorname{Fr}_r^f
\]
the inner-diagonal group is the adjoint fixed-point group.  The pinned
part is generated by $\operatorname{Fr}_r$ and the diagram centralizer
of $\gamma$, with the relation
$\gamma\operatorname{Fr}_r^f=1$ on fixed points.  In the ordinary
field--Dynkin presentation the restriction kernel is exactly the
cyclic group generated by $F_0$; see \cite[\S2.2]{SpaethD}.  These operations act
on the partial regular extension, including its added torus, with the
same relation.  The exceptional graph-isogeny cases have instead a
generator $\tau$, with $\tau^2$ a split Frobenius, and defining
endomorphism $F_0=\tau^f$ \cite[Definition~3.1]{BMO}.  Its restriction
has order exactly $f$.  For untwisted groups, $f=2d$, even powers are
field operations of order $d$, and odd powers exchange long and short
root groups.  For the Suzuki and Ree twists, $f$ is odd and $\tau^2$
has field order $f$, forcing $\tau$ itself to have order $f$.
The usual automorphism description is
faithful on the universal component after the finite exceptional list
has been removed.  We now turn its field operations into actual
algebraic automorphisms; no field Frobenius is being regarded as an
invertible morphism of algebraic groups.

First consider the stabilizer of one component in the pinned image of
$S$.  In the ordinary case its projection modulo diagrams to the field
cyclic group has order $a$, where $a\mid f$ is a $p$-power and
$a\leq |S|$.  Put $d=f/a$.  On $a$ copies of the pinned group define
\begin{align*}
 R(g_0,\ldots,g_{a-1})
    &=(\gamma(g_{a-1}),g_0,\ldots,g_{a-2}),\\
 F_a&=R\operatorname{Fr}_r^d.
\end{align*}
The Frobenius in the second formula acts on every coordinate.  We have
$R^a=\gamma$ diagonally.  Projection onto coordinate zero identifies
$(\mathbf J_0^a)^{F_a}$ with $\mathbf J_0^{F_0}$: a fixed tuple is
\[
 (g,\operatorname{Fr}_r^d(g),\ldots,
       \operatorname{Fr}_r^{(a-1)d}(g)),
 \qquad F_0(g)=g.
\]
On these tuples $R$ induces $\operatorname{Fr}_r^{-d}$.
Every diagram $\delta$ centralizing $\gamma$ acts diagonally and
commutes with $R$.  Thus the subgroup above the order-$a$ field image
is represented with its exact presentation
\[
 R^a=\gamma,\qquad [R,\delta]=1,
\]
as well as the relations among the diagrams.  The required pinned
$p$-subgroup is the corresponding subgroup of this finite group.
One does not replace $R$ by an order-$a$ permutation when
$\gamma\ne1$: the displayed wrap relation is essential.  In particular
mixed field--graph elements and their relations are retained.

For $F_0=\tau^f$, the pinned automorphism group is cyclic.  If its
relevant subgroup has order $a$, use $a$ copies, the pure cyclic
permutation $R$, and
\[
 F_a=R\tau^{f/a}.
\]
Projection identifies the fixed points with those of $\tau^f$, and
$R$ induces $\tau^{-f/a}$.  The permutation is an algebraic
automorphism even though $\tau$ need not be one.  Both constructions
have a power of $F_a$ equal to a Frobenius endomorphism.

For completeness, the passage from a component stabilizer to its
orbit preserves relations.  Choose a base component and, for every
component in its orbit, an element of the pinned image transporting
the base component to it.  Use these transports to choose the
identifications.  If an element sends component $i$ to component $j$,
its map in these identifications is the stabilizer element obtained
by composing the inverse chosen transport to $j$, that element, and
the chosen transport to $i$.  Products of these maps satisfy the
stabilizer multiplication law, by cancellation of the middle
transports.  Thus they act on the product of the copy models by
permutations and the algebraic stabilizer operations just constructed.
The resulting group is the actual pinned image of $S$, and is a
$p$-group.  The orbit size and all copy numbers are bounded by $|S|$.

Apply the same construction to the adjoint groups.  Their fixed points
identify with the original inner-diagonal groups, and the action of
the pinned image agrees with its original action.  The given
representatives consequently define elements of the asserted
semidirect product.  The standard faithful automorphism description
ensures that a relation on the universal components has exactly its
indicated inner image in this model.  This also applies to the actions
extended to the torus in the partial regular extension: the pinned
relations hold there by construction.  Hence the representative
product $h_xh_y=n_{xy}h_{xy}$ still has that inner factor, rather than
merely an unspecified inner automorphism.
\end{proof}

Write $b'$ for the resulting product block of $J$; it is $I$-stable.
The number of original components and $|S|$ are bounded by
Lemma~\ref{lem:ext:reduced-structure}.  Thus the number of absolute
components in the new model is bounded, although their classical
ranks and the central torus rank need not be bounded.

\subsection{A central twist with an invariant parabolic}

Let $s\in\mathbf J^{*F}$ be a semisimple $p'$-element such that $b'$
belongs to $\mathcal E_p(J,s)$.  The partial regular extension ensures
that
\[
 \mathbf C=C_{\mathbf J^*}(s)
\]
is connected.  Indeed its possible semisimple-centralizer component
obstruction is $p$-primary, since the component group of
$Z(\mathbf J)$ is $p$-primary; the order of a component arising from
$s$ also divides the order of $s$.  These orders are coprime.  This is
the semisimple-centralizer component criterion, applied to the dual
root datum; see \cite[Corollary~2.9]{BonnafeQuasiIsolated} and
\cite[Proposition~14.20]{MalleTesterman}.  In the exceptional-isogeny cases the obstruction has no
nontrivial points, as in Lemma~\ref{lem:ext:partial-extension}.

\begin{lemma}[The fixed central torus]\label{per:fixed-torus}
There exist an $F$-stable torus $\mathbf V\subseteq Z^\circ(\mathbf C)$,
a dual Levi subgroup
\[
 \mathbf L^*=C_{\mathbf J^*}(\mathbf V),
\]
and a bounded positive integer $m_0$ such that
\[
 \mathbf C\subseteq\mathbf L^*,\qquad
 z_m=s^{p^m-1}\in\mathbf V
       \quad\hbox{whenever }m_0\mid m.
\]
The $s$-transporters of the pinned image fix $\mathbf V$ pointwise.
Each $z_m$ defines canonically a linear character $\lambda_m$ of
$L=\mathbf L^F$ of $p'$-order, through the quotient torus
$\mathbf L/[\mathbf L,\mathbf L]$.
\end{lemma}

\begin{proof}
Diagonal automorphisms do not change a geometric series parameter.
Since $b'$ is stable, each element of $Y'$ preserves the geometric
class of $s$.  Connectedness of $\mathbf C$ and Lang's theorem imply
that this geometric class contains a unique rational class.  Hence
\[
 H_s=\{(g,y)\in\mathbf J^{*F}\rtimes Y':
                         g\,y(s)g^{-1}=s\}
       \longrightarrow Y'
\]
is onto, with kernel $\mathbf C^F$.  That kernel acts trivially on
$\mathbf V_0=Z^\circ(\mathbf C)$, so $H_s$ defines an action of $Y'$
on $\mathbf V_0$.  This action commutes with $F$.  Set
\[
 \mathbf V=(\mathbf V_0^{Y'})^\circ.
\]
In particular the action just defined is independent of all choices
of transporters.

We show that a bounded $p'$-power of $s$ lies in $\mathbf V_0$.
Choose a maximal torus of $\mathbf C$, and let $X$ be its character
lattice, $Q$ the ambient root lattice, and $Q_s$ the lattice generated
by roots centralizing $s$.  Since the primal derived group is simply
connected, $Q$ is saturated in $X$.  The torsion of $X/Q_s$ therefore
injects into the torsion of $Q/Q_s$.  In type A the subsystem is a
product of equality packets, and its lattice is primitive in the
ambient root lattice.  In the classical signed-coordinate systems,
orient each inverse-paired eigenvalue packet.  A type-A packet then
has its integral zero-sum lattice.  On each residual B, C, or D
packet, the lattice generated by the centralizing roots contains
twice every integral coordinate vector in its rational span.
Consequently the saturation quotient on each such packet has
exponent dividing two; taking products does not enlarge this
exponent.  This proves a uniform exponent bound for classical
components.  On the exceptional components there are only finitely
many root subsystems of the finitely many bounded root systems, so
their saturation quotients also have bounded exponent.  The bounded
number of absolute components allows one common exponent.

The character group of the component group of $Z(\mathbf C)$ is
the torsion just considered.  Since $s$ is central in $\mathbf C$,
there is therefore a bounded integer $l$, prime to $p$, with
$s^l\in\mathbf V_0$.  We may discard the $p$-part of the exponent
because $s$ has $p'$-order.  This element is fixed by $Y'$.
For any torus with a $Y'$-action, the norm morphism
\[
 t\longmapsto\prod_{y\in Y'}y(t)
\]
has connected image in the fixed torus.  On a fixed element it is
the $|Y'|$th power.  Thus the group of components of the fixed torus
on geometric points is killed by $|Y'|$, a $p$-power.  Its element
represented by $s^l$ has $p'$-order, and is trivial.  Hence
$s^l\in\mathbf V$.

Take $m_0$ with $p^{m_0}\equiv1\pmod l$; for $l=1$ take $m_0=1$.
This is bounded, and proves the assertion about $z_m$.  The torus
$\mathbf V$ is $F$-stable, and its centralizer is an $F$-stable Levi.
Since $\mathbf V\subseteq Z(\mathbf C)$, the Levi $\mathbf L^*$
contains $\mathbf C$.
Moreover $\mathbf V\subseteq Z^\circ(\mathbf L^*)$.
The natural duality
\[
 Z^\circ(\mathbf L^*)
    \quad\longleftrightarrow\quad
 \mathbf L/[\mathbf L,\mathbf L]
\]
associates to its rational point $z_m$ a linear character
$\lambda_m$ of $L$, trivial on the finite points of the algebraic
derived subgroup.  It has $p'$-order and takes values in $\cO$.
This construction is functorial for the transporter operations.
\end{proof}

\begin{lemma}[Common parabolic representatives]\label{per:parabolic}
The primal Levi $\mathbf L$ can be placed in a parabolic $\mathbf Q$
such that, for the subgroup $A_1$ of $I$ preserving $\mathbf Q$,
$\mathbf L$, the rational $s$-series in $L$, and $\lambda_m$, one has
\begin{equation}\label{per:parabolic-factorization}
 I=\Inn(J)A_1,\qquad A_1\cap\Inn(J)=\Inn(L).
\end{equation}
Here inner groups mean their images in the algebraic automorphism
group.  Only $\mathbf L$ is required to be $F$-stable.
\end{lemma}

\begin{proof}
Choose a cocharacter $\nu$ of $\mathbf V$ outside the finitely many
root hyperplanes which do not contain all of $\mathbf V$.  Then
$C_{\mathbf J^*}(\nu)=\mathbf L^*$, and $\nu$ specifies a parabolic
$\mathbf Q^*$ with this Levi.  Every element of $H_s$ fixes
$\mathbf V$ pointwise, and hence preserves both this Levi and the
positive root set of this parabolic.  There is no requirement that
$\nu$ be $F$-fixed.

Choose an $F$-stable Borel--torus pair in the connected group
$\mathbf C$, and denote its torus by $\mathbf T^*$.  The Lang--Steinberg theorem
\cite[Theorem~21.7]{MalleTesterman} makes the rational Borel--torus pairs a single $\mathbf C^F$-orbit.
Thus each $s$-transporter can be corrected by an element of
$\mathbf C^F$ to normalize $\mathbf T^*$.  Such a correction still
fixes $\mathbf V$ pointwise.  On this torus the transporter is a
Weyl operation followed by its specified pinned operation, and
commutes with the rational root datum endomorphism.

Choose the matching rational primal torus $\mathbf T$.  Root--coroot
duality transfers the zero and positive root sets to a Levi and a
parabolic $\mathbf L\subseteq\mathbf Q$ of $\mathbf J$.  Equivalently,
an invariant rational positive definite form transfers the
inequalities defined by $\nu$ to primal coroot inequalities; a
multiple clears denominators.  The zero set is $F$-stable and is
exactly the dual Levi root system.  The transferred transporter
preserves these positive and zero sets.

We spell out why its primal representative can be rational.  Its
Weyl and pinned operation has a geometric representative in
$\mathbf J_{\mathrm{ad}}\rtimes Y'$.  Compatibility with the rational
torus endomorphism says that its discrepancy with $F$ lies in
$\mathbf T_{\mathrm{ad}}$.  Lang's theorem for this connected torus
corrects that discrepancy without changing the root action or its
positive set.  This gives an element of
$\mathbf J_{\mathrm{ad}}^F\rtimes Y'$ preserving $\mathbf Q,\mathbf L$
and the indicated torus data.  In particular it preserves the
rational $s$-class and $\lambda_m$.

One may also prescribe its diagonal coset.  For any rational
maximal torus of $\mathbf J$, comparison of the exact sequences for
\[
 Z(\mathbf J)\longrightarrow\mathbf T\longrightarrow
 \mathbf T_{\mathrm{ad}},\qquad
 Z(\mathbf J)\longrightarrow\mathbf J\longrightarrow
 \mathbf J_{\mathrm{ad}}
\]
and Lang's theorem for $\mathbf T$ and $\mathbf J$ give a surjection
\[
 \mathbf T_{\mathrm{ad}}^F\longrightarrow
       \mathbf J_{\mathrm{ad}}^F/\operatorname{im}(J).
\]
Multiplying the representative by such a torus element changes its
diagonal coset arbitrarily and preserves the parabolic and Levi.
Diagonal operations preserve geometric series and quotient-torus
characters.  Moreover $C_{\mathbf L^*}(s)=\mathbf C$ is connected,
so preservation of the geometric class here is preservation of its
unique rational class.  Given an element of $I$, choose the diagonal
coset to agree with that element.  The resulting representative
differs from it by an element of $\Inn(J)$, and therefore belongs to
$I$.  This proves the first equality.

Finally, the simultaneous normalizer of a parabolic and its chosen
Levi in the connected group is the Levi itself.  Therefore an
inner operation of $J$ lies in $A_1$ exactly when its implementer
lies in $L$.  Elements of $L$ also preserve its $s$-class and linear
character.  This proves the second equality.
\end{proof}

\subsection{The integral comparison and its ordinary characters}

The chosen Levi contains the full dual centralizer, and its
parabolic admits the required automorphism representatives.
These are the geometric inputs for an equivariant integral
Morita equivalence.  The remaining character calculation will
return the Levi block after a bounded Frobenius power and a
central linear twist.

\begin{lemma}[Equivariant integral Levi equivalence]\label{per:br}
There is an $A_1$-stable block $b_L$ of $\cO L$ and an integral Morita
bimodule $U$ between $\cO Jb'$ and $\cO Lb_L$ with a strict
$A_1$-action.  For $\ell\in L$, the operator of
$\operatorname{ad}(\ell)\in A_1$ on $U$ is exactly
\begin{equation}\label{per:br-inner}
 u\longmapsto\ell u\ell^{-1}.
\end{equation}
Every central element of $J$ acts identically from the two sides of
$U$.
\end{lemma}

\begin{proof}
We use the full-centralizer form of the integral
Bonnaf\'e--Rouquier theorem \cite[Theorem~B$'$, \S11.4]{BR}, also
stated in \cite[\S7]{BDR}.  Its hypotheses here are that
$\mathbf J$ is connected reductive in characteristic $r\ne p$,
$F$ has a Frobenius power, $\mathbf L$ is an $F$-stable Levi of a
parabolic $\mathbf Q$, $s$ is a rational semisimple $p'$-element,
and $C_{\mathbf J^*}(s)\subseteq\mathbf L^*$.  An $F$-stable
parabolic is not required.  All these hypotheses were established
above.  With $\mathbf Q=\mathbf L\mathbf U_Q$, the theorem gives
the Morita bimodule in the single nonzero degree of the appropriate
block cut of compactly supported cohomology of
\[
 Y_{\mathbf U_Q}=
 \{g\mathbf U_Q:g^{-1}F(g)\in
                         \mathbf U_QF(\mathbf U_Q)\}.
\]
Its degree is
$\dim\mathbf U_Q-\dim(\mathbf U_Q\cap F(\mathbf U_Q))$.

The construction is valid over the fixed unramified ring $\cO$.
Indeed the group-algebra block idempotents are defined over some
finite unramified Witt subring: lift their finite-residue-field
idempotents uniquely.  Construct the compactly supported integral
cohomology complex and these cuts over that ring.  After faithfully
flat extension to a splitting discrete valuation ring the cited
integral theorem gives concentration, projectivity on both sides,
and the Morita evaluation isomorphisms.  These properties descend
by faithful flatness.  Extending the descended construction to
$W(k)$ gives $U$.  Thus extension to splitting coefficients is used
to test the equivalence, not as an assumption that the generic field
of $W(k)$ contains $p$-power roots of unity.

For $a\in A_1$ the map $g\mathbf U_Q\mapsto a(g)\mathbf U_Q$ is
an actual automorphism of this variety.  Using inverse pullback
consistently gives a strict action on cohomology.  The series cuts
are preserved; since $b'$ is stable, equivariance of the block
correspondence makes $b_L$ stable too.  For $a=\operatorname{ad}(\ell)$,
$\ell\in L$, this variety map is precisely the left action of
$\ell$ followed by the right action of $\ell^{-1}$.  This proves the
asserted identity.  A central element of $J$, which lies in $L$,
gives the identity variety map, proving central compatibility.
\end{proof}

\begin{lemma}[Returning the Levi block]\label{per:row-return}
After replacing $m_0$ by a bounded multiple $m$, one has
\begin{equation}\label{per:levi-return}
 \sigma^m(b_L)=\lambda_m b_L,
\end{equation}
in the convention that multiplication of ordinary characters by
$\lambda_m$ gives the block on the right.  The character $\lambda_m$
is still defined by $s^{p^m-1}$ and is $A_1$-invariant.
\end{lemma}

\begin{proof}
Define an operation on the ordinary characters of $L$ by
\[
 T(\chi)=\lambda_{m_0}^{-1}\sigma^{m_0}(\chi).
\]
Coefficient Frobenius and twisting by a linear character both
permute ordinary irreducible characters and their blocks.  It
therefore suffices to find a bounded positive integer $d$ such
that $T^d$ fixes one ordinary irreducible character in $b_L$: its
block then returns as well.  We first bound row orbits after a
regular embedding, then use restriction to return a row of $b_L$.
At the end we check that the accumulated twist is precisely
$\lambda_{d m_0}$.

For an ordinary character in $\mathcal E(L,st)$, with $t$ a
$p$-element of $C_{\mathbf L^*}(s)^F$, coefficient Frobenius sends
the torus parameter to $s^{p^{m_0}}t$.  The $p$-part is unchanged
by our choice of the characteristic-zero extension of $\sigma$.
The Deligne--Lusztig character formula, whose Green functions are
rational, gives the same statement for every uniform pairing.
Thus $T$ returns the parameter to $st$ and fixes its uniform pairings.
This argument concerns actual pairings, not just the set of torus
characters; compare \cite[Lemma~3.6(i)]{FarrellKessar}.

\emph{Bounded row orbits after regular embedding.}
Use a full regular embedding
$\mathbf L\hookrightarrow\mathbf L^\dagger$ with connected center
and compatible Steinberg endomorphism.  The derived subgroup stays
simply connected.  The standard regular embedding construction
applies also to Frobenius roots; in the exceptional-isogeny cases
at issue the centers on points already have no obstruction.
On dual groups the map
\[
 \pi:\mathbf L^{\dagger *}\longrightarrow\mathbf L^*
\]
has central torus kernel.  Choose a character above a row of $b_L$,
with parameter $s^\dagger t^\dagger$, where $s^\dagger$ is its
$p'$-part and $t^\dagger$ its $p$-part.  After rational conjugacy its
$p'$-parameter projects to $s$.  This rational adjustment is possible
by connectedness of the centralizer and Lang's theorem.  Put
\[
 z^\dagger=(s^\dagger)^{p^{m_0}-1}.
\]
Every root of $\mathbf L^{\dagger *}$ has value one on
$z^\dagger$, by the root test already established on $\mathbf L^*$.
Thus $z^\dagger$ is central.  The center of
$\mathbf L^{\dagger *}$ is connected because
$[\mathbf L^\dagger,\mathbf L^\dagger]$ is simply connected.
The associated linear character $\lambda^\dagger$ restricts to
$\lambda_{m_0}$.  The operation
\[
 T^\dagger(\chi)=(\lambda^\dagger)^{-1}\sigma^{m_0}(\chi)
\]
therefore preserves $\mathcal E(L^\dagger,s^\dagger t^\dagger)$
and fixes each of its uniform pairings.

Ordinary Jordan decomposition on these connected data identifies
the pairings with the unipotent uniform pairings of the connected
centralizer.  Separate the ordinary classical factors from the
bounded-rank exceptional factors.  The latter include every factor
with an exceptional graph-isogeny endomorphism, in particular the
Suzuki type \({}^2B_2\), and every triality factor.  On the ordinary
classical factors the pairings distinguish the irreducible unipotent
characters, even up to scalar: type A is uniform, and for the ordinary
types B, C, and D this is the Fourier separation of the symbol labels
in Lemma~\ref{lab:fingerprints}.  That separation concerns ordinary
characters and is independent of the coefficient prime $p$.
Every uniform pairing vector here is nonzero.  Varying the torus
test on one factor while fixing nonzero tests on the others shows
that equality of product pairing vectors makes the corresponding
factor vectors proportional.  Lemma~\ref{lab:fingerprints} then
fixes all ordinary classical labels.  Equality uniqueness for a
centralizer consisting of such factors is also the ordinary character
statement of \cite[Lemma~3.7]{FarrellKessar}.
Consequently there is at most one choice on each ordinary classical
factor.  Every remaining factor has bounded root system and one of
finitely many rational twists, so its number of unipotent labels is
bounded.  There are boundedly many such factors.  Adding central tori
introduces no new unipotent
labels.  It follows that the orbit of $\chi$ under $T^\dagger$ has
bounded length, independently of the ranks of the classical factors
and of the field cardinalities.  Frobenius-permuted factors are
calculated on a representative with the composite endomorphism;
the same reasoning applies to them.

\emph{Restriction to the original Levi.}
The upstairs operation $T^\dagger$ has bounded row orbits,
uniformly in rank and field size.  For the chosen
$\lambda^\dagger$, restriction satisfies
\[
 \Res_L\bigl(T^\dagger(\chi)\bigr)
       =T\bigl(\Res_L\chi\bigr),
\]
because $\lambda^\dagger|_L=\lambda_{m_0}$ and coefficient
Frobenius commutes with restriction.  If $(T^\dagger)^a$ fixes an
upstairs row, then $T^a$ permutes its restriction constituents.  We now
distinguish the type of the original ambient components, rather
than the centralizer factors used in the pairing argument.  We
are working on one $S$-orbit, so these original components have
the same type.

For a non-type-A orbit the number of distinct restriction
constituents is bounded.  For a common maximal torus of the primal
groups $\mathbf L\subseteq\mathbf J$, let $X$ be its character
lattice and $\Phi_L\subseteq\Phi_J$ their root systems.  The torsion in
$X/\mathbb Z\Phi_L$ injects into the torsion in
$X/\mathbb Z\Phi_J$: after making the Levi standard, the kernel
$\mathbb Z\Phi_J/\mathbb Z\Phi_L$ is free on the omitted simple
roots.  The center component group of the Levi therefore has order
bounded by that of the original non-A group, with a bounded product
over the copies.  The rational quotient
\[
 (\mathbf L^\dagger)^F/
       \bigl(Z(\mathbf L^\dagger)^F L\bigr)
\]
has bounded order by the central exact sequence and Lang's theorem.
Clifford theory bounds the number of restriction constituents by
this order.  A bounded power of $T^\dagger$ fixes the upstairs row,
so the same power of $T$ permutes only this bounded set of
downstairs constituents.  A further bounded power fixes a
constituent belonging to $b_L$.

For a type-A orbit this center-component bound need not hold.
Choose a covering block of $b_L$ in $L^\dagger$ and apply the
preceding upstairs construction to a row
\[
 \chi^\dagger\in\mathcal E(L^\dagger,s^\dagger)
\]
in that block.  Such a row exists by the integral basic-set theorem
for type A with connected center: intersect its basic set with the
covering block.  Its parameter projects to the prescribed $s$.
The resulting $\lambda^\dagger$ still restricts to $\lambda_{m_0}$,
so the downstairs operation $T$ is unchanged.
We claim that $\Res_L\chi^\dagger$ is irreducible.  The preimage
\[
 \pi^{-1}(C_{\mathbf L^*}(s))
\]
is connected, since both its kernel and its quotient are connected.
It centralizes $s^\dagger$: a maximal torus and its root subgroups
generate it, and their roots take value one on $s^\dagger$ exactly
when they do on $s$.  Hence if $s^\dagger u$ is conjugate to
$s^\dagger$, with $u\in\ker\pi$, then any conjugating element is
in this preimage and $u=1$.

The quotient linear characters of $L^\dagger/L$ correspond
faithfully to rational points of the dual kernel and act by these
central multiplications on parameters.  Only the trivial quotient
character can therefore fix $\chi^\dagger$.  Frobenius reciprocity
and induction from a normal subgroup with abelian quotient give
\[
 \bigl\langle\Res_L\chi^\dagger,
                    \Res_L\chi^\dagger\bigr\rangle
 =\bigl|\{\eta\in\Irr(L^\dagger/L):
                         \eta\chi^\dagger=\chi^\dagger\}\bigr|=1.
\]
This proves the claim; equivalently one may use
\cite[Lemma~4.8]{SFTV}.  A covering block covers a single orbit of
blocks of $L$.  Since this restriction is irreducible and
$L^\dagger$-invariant, its block is the prescribed covered block
$b_L$.  Its orbit under $T$ is now
bounded by the upstairs orbit, with no diagonal-index bound.

In either case $T^d$ fixes a row of $b_L$ for a bounded positive
integer $d$.  Iteration has exactly the required twist prescription:
\[
 \prod_{i=0}^{d-1}\sigma^{i m_0}(\lambda_{m_0})
       =\lambda_{d m_0},\qquad
 (p^{m_0}-1)\sum_{i=0}^{d-1}p^{i m_0}=p^{d m_0}-1.
\]
Indeed the first identity follows from torus duality and the second.
Thus $\lambda_{d m_0}^{-1}\sigma^{d m_0}$ fixes that row and hence
its block.  Taking this bounded multiple proves the assertion.
The construction of Lemma~\ref{per:fixed-torus} continues to make
$\lambda_m$ invariant under $A_1$.
\end{proof}

\subsection{Scalar overlap and the specified inner factors}

\begin{lemma}[Projective extension of the transports]
\label{per:scalar-overlap}
For the Lie-component orbit above there is an integral Morita
bimodule $\mathcal M_J$ between $\sigma^m(b')$ and $b'$ on which
$I$ has covariance transports defined up to scalars in $\cO^\times$.
For every $g\in J$ its inner transport agrees, modulo scalars, with
\[
 E_g(v)=gvg^{-1}.
\]
Central $p$-elements act equally from both sides of $\mathcal M_J$.
\end{lemma}

\begin{proof}
Put $A=\cO Jb'$, $B=\cO Lb_L$, and use primes for their
$\sigma^m$-images.  By Lemma~\ref{per:row-return}, multiplication
of characters by $\lambda_m$ identifies $B$ with $B'$.
Let $T$ be the associated invertible $(B',B)$-bimodule.
Explicitly its underlying left module is $B'$, and its right action
uses the isomorphism
\[
 f:B\longrightarrow B',\qquad
        g b_L\longmapsto\lambda_m(g)^{-1}g\sigma^m(b_L).
\]
Let $U^\vee$ be a Morita inverse of the bimodule in
Lemma~\ref{per:br}.  Then
\[
 \mathcal M_J=
   \sigma^m(U)\otimes_{B'}T\otimes_B U^\vee
\]
is an $(A',A)$-Morita bimodule.  Each factor has a strict $A_1$-action:
on the outer factors this is the geometric action and its dual;
on $T$ it is the action on its group basis, since $\lambda_m$ is
$A_1$-invariant.  Hence their tensor product has a strict action,
which we denote by $D_a$.

The strict action of $A_1$ has been constructed.  To extend it
projectively to $I=\Inn(J)A_1$, we must compare it with the
actual inner operators on their common subgroup.
The comparison on the intersection in
Lemma~\ref{per:parabolic} is explicit.  If $\ell\in L$, the outer
factors identify the geometric action with the actual left and
inverse right action.  On the middle factor the right action uses
$f(\ell^{-1})=\lambda_m(\ell)\ell^{-1}$, so
\begin{equation}\label{per:overlap-scalar}
 D_{\operatorname{ad}(\ell)}
          =\lambda_m(\ell)^{-1}E_\ell.
\end{equation}
All discrepancies therefore lie in the scalar units; no unit in
the non-scalar part of the block center is introduced.

The actual inner operators satisfy
$E_gE_h=E_{gh}$ and
$D_aE_gD_a^{-1}=E_{a(g)}$.  By
$I=\Inn(J)A_1$, represent an element of $I$ as
$\operatorname{ad}(g)a$ and use $E_gD_a$.  Any two such
representations differ in the intersection $\Inn(L)$, where the
preceding formula shows that the two operators differ by a scalar.
A central ambiguity in $g$ is also scalar: a central $p'$-element
acts by its central character on each block, while a central
$p$-element gives no discrepancy.  For the latter assertion apply
the displayed overlap formula to a central $p$-element $c$.  Its
algebraic inner operation is the identity and
$\lambda_m(c)=1$, whence $E_c=\id$.  Covariance then shows that
multiplication of the chosen operators agrees with multiplication
in $I$ modulo scalars, as asserted.
\end{proof}

The adjective scalar in Lemma~\ref{per:scalar-overlap} is necessary.
The following elementary observation identifies exactly how it is
used, and records the role of the element-wise factor identity.

\begin{lemma}[Removing the scalar defect]\label{per:scalar-cocycle}
Let $S$ be a finite $p$-group acting on two block algebras by crossed
systems with the same group labels and respective factors
$a_{xy},a'_{xy}$.  Suppose that a nonzero Morita bimodule has
invertible covariance maps $D_x$ such that
\[
 D_xD_y(v)=c(x,y)a'_{xy}D_{xy}(v)a_{xy}^{-1},
       \qquad c(x,y)\in k^\times.
\]
If the factors obey their actual crossed-system identities, the
maps can be rescaled to satisfy \eqref{per:factor-law} exactly.
\end{lemma}

\begin{proof}
Normalize $D_1=\id$.  Compare $(D_xD_y)D_z$ and $D_x(D_yD_z)$.
Covariance and
\[
 a_{xy}a_{xy,z}=\alpha_x(a_{yz})a_{x,yz},
 \qquad
 a'_{xy}a'_{xy,z}=\alpha'_x(a'_{yz})a'_{x,yz}
\]
cancel all the non-scalar factors.  Because the $D_x$ are
$k$-linear, the remaining equality is
\[
 c(x,y)c(xy,z)=c(y,z)c(x,yz).
\]
Thus $c$ is an ordinary scalar $2$-cocycle, with trivial action on
$k^\times$.  Positive-degree cohomology of $S$ is annihilated by
$|S|$.  The $|S|$th-power map on $k^\times$ is an automorphism,
since $k$ is algebraically closed of characteristic $p$.
It induces an automorphism on this cohomology, which must therefore
vanish.  In particular $H^2(S,k^\times)=0$.  Rescale $D_x$ by a
scalar $1$-cochain trivializing $c$ to obtain the required maps.
\end{proof}

\subsection{Tensor products and central quotients}

The Lie-component comparison now has only scalar multiplication
errors.  We assemble it with the remaining component families,
descend through the original central kernel, and only then remove
the scalar error over $k$.  Descending before invoking integral
finiteness is essential because the lifted blocks may have
unbounded defect.

\begin{proof}[Proof of Proposition~\ref{prop:sylow-comparison}]
On each cross-characteristic Lie orbit use
Lemma~\ref{per:scalar-overlap}.  Choose a common bounded multiple
of its comparison exponents.  Passing to a multiple is legitimate:
tensor successive Frobenius twists of the comparison bimodule and
of its transports.  The twist characters multiply according to the
identity in Lemma~\ref{per:row-return}; scalar projectivity and the
equality of central $p$-actions are preserved.  The number of
orbits is bounded.  Thus this choice still gives a bounded positive
integer $m$.

The other universal components are handled directly.  For the
finite list of groups, take a common Frobenius period of their block
idempotents and use the identity bimodule, with the ordinary group
operations as transports.  Permuted isomorphic components use
matching models, and permutation of tensor factors gives their
transports.  The finite list includes the bounded exceptional
multipliers and automorphisms already removed from the Lie models.

For the unbounded alternating family the universal cover is
$2.\mathfrak A_n$, outside a finite list.  Every automorphism is
induced by $2.\mathfrak S_n$: use the automorphisms of
$\mathfrak A_n$ and unique lifting to its universal cover
\cite[Theorem~2.3 and \S\S2.7.2--2.7.3]{Wilson}.  We give a uniform period that also
handles the double cover.  Let $u$ be an integer coprime to the
exponent of $2.\mathfrak A_n$.  For $g\in2.\mathfrak A_n$,
its image and the image of $g^u$ have the same cycle type in
$\mathfrak S_n$.  For a suitable $t\in2.\mathfrak S_n$ and the
central involution $z$ this gives
\[
 g^u=z^\epsilon t g t^{-1}.
\]
Since $u$ is odd,
\[
 g^{u^2}
   =z^{\epsilon(u+1)}t^2g t^{-2}
   =t^2g t^{-2},\qquad t^2\in2.\mathfrak A_n.
\]
Thus every square of a cyclotomic powering fixes every conjugacy
class.  The cyclotomic action of coefficient Frobenius is powering
by such a unit $u$ (chosen to be $p$ on the $p'$-part and $1$ on the
$p$-part of the exponent).  Its square fixes all ordinary
characters and hence all block idempotents.  After making $m$ even,
the identity block bimodule and its ordinary automorphism
transports apply also to the alternating orbits.  Blocks inflated
from the simple quotient cause no change in this argument.

On the factor $P$ take the identity bimodule for $\cO P$, with the
given automorphism transports.  On the central $p'$-group $Z$ a
block is the rank-one character algebra for some character $\theta$.
Identify it with $\cO$ and compare it with the rank-one algebra for
$\sigma^m\theta$.  The representatives act trivially on $Z$, so
use the identity transport; inner left/right discrepancies on this
factor are scalars.  These choices are compatible with permutation
of the other tensor factors.  On Lie orbits compatibility was built
into a single product algebraic group and its actual variety
actions, so no coherence choice between individual components is
needed.

Tensor all these bimodules over $\cO$.  We now descend from the
product cover, including its added direct central $p$-factors, to
$\bar N$.  Here is the precise quotient argument.  If a central
$p$-element $c$ acts equally from the two sides of a Morita bimodule
$V$, then
\[
 (c-1)V=V(c-1).
\]
For the two-sided ideals generated by any collection of such
elements, the Morita correspondence consequently matches the
ideals.  The quotient of $V$ by their common action is a Morita
bimodule between the quotient algebras; this also follows by
quotienting the Morita evaluation isomorphisms.  Apply this to the
added direct central $p$-factors and to the $p$-part of the old
central kernel $D_0$.  Every factor construction has the required
equality, so the tensor product does too.  For the $p'$-part of
$D_0$, the chosen block lifts have trivial kernel character on both
sides, including after Frobenius.  This part of the kernel already
acts trivially.  No bound for $|D_0|$ is required.  The quotient is
therefore the claimed integral Morita bimodule $\mathcal M$.

We now have the integral comparison on the original quotient
group, where the defect bounds apply.  It remains to recover the
specified representative factors and normalize their scalar error
after reduction.
The covariance operators preserve the quotient ideals, and descend
with it.  On the product cover, products of the operators for $h_x$
and $h_y$ differ from that for $h_{xy}$ by the actual inner operator
of a lift of $n_{xy}$, up to a scalar.  This follows on each Lie
orbit from Lemma~\ref{per:algebraic-action} and
Lemma~\ref{per:scalar-overlap}, and on the remaining factors from
their explicit ordinary transports.  Choices of a lift differ by
$D_0$; after descent their inner operators agree.  Consequently on
$\mathcal M$ the product law is the specified law for $n_{xy}$ up
to a scalar.  The factors in $\bar N$ satisfy the actual equality
\[
 n_{xy}n_{xy,z}=\alpha_x(n_{yz})n_{x,yz},
\]
by their construction from the chosen representatives.  Reduce
modulo $p$ and apply Lemma~\ref{per:scalar-cocycle}.  The resulting
normalized maps are exactly \eqref{per:covariance} and
\eqref{per:factor-law}.

All exponents used above depend only on the bounded number of
components and copies, the order of $S$, the bounded exceptional
root systems, and the finite exceptional group list.  These data
are bounded in terms of $p,M$ by the extension reductions.
The classical root-lattice exponent, the classical row separation,
and the alternating powering argument were uniform in the ranks
and field sizes.  This proves the asserted uniformity and the
proposition, including its trivial-subgroup case.
\end{proof}

\subsection{Extension of the coefficient field}
\label{sec:coefficient-extension}

The preceding argument works over $k=\overline{\F}_p$.  We finish by
transferring its finite list to an arbitrary algebraically closed field
of characteristic $p$.  The block comparison is the coefficient-field
argument of \cite[Section~2.1]{BlockwiseAlperinWeight}; we recall the
part needed here and combine it with scalar extension of Morita
bimodules.

\begin{lemma}[Coefficient extension]
\label{lem:coefficient-extension}
Let $k=\overline{\F}_p$, let $K$ be an algebraically closed field of
characteristic $p$, and choose an embedding $k\hookrightarrow K$.
For every finite group $G$, the map $kG\longrightarrow KG$ induces a
bijection on block idempotents, sending $b$ to $b_K=1\otimes b$.
The blocks $kGb$ and $KGb_K=K\otimes_k kGb$ have the same defect
groups as $p$-subgroups of $G$.

If finite-dimensional $k$-algebras $A$ and $A'$ are $k$-linearly
Morita equivalent, then $K\otimes_k A$ and $K\otimes_k A'$ are
$K$-linearly Morita equivalent.
\end{lemma}

\begin{proof}
The center commutes with scalar extension:
\[
 K\otimes_k Z(kG)\simeq Z(KG).
\]
Indeed, the center is the kernel of the map sending an element to its
commutators with the finitely many elements of $G$, and field extension
preserves kernels.  Write $Z(kG)=\prod_i Z_i$ as a product of local
finite-dimensional commutative $k$-algebras.  Since $k$ is algebraically
closed, each $Z_i$ has residue field $k$ and nilpotent radical.  In
$K\otimes_k Z_i$ the extended radical is nilpotent and the quotient is
$K$, so this algebra is again local.  Thus the primitive central
idempotents correspond bijectively.

For a $p$-subgroup $P\le G$, the Brauer map deletes the coefficients
outside $C_G(P)$, and hence
\[
 \Br_P^K(1\otimes b)=1\otimes\Br_P^k(b).
\]
The map $kC_G(P)\longrightarrow KC_G(P)$ is injective.  The two Brauer
images are therefore nonzero for exactly the same $P$.  Their maximal
such $p$-subgroups, which are the defect groups by the convention in
Section~\ref{sec:preliminaries}, are identical.

For the Morita assertion, choose finite-dimensional inverse Morita
bimodules ${}_{A'}U_A$ and ${}_A V_{A'}$.  The $k$-linearity of the
equivalence means that the left and right $k$-actions on these bimodules
agree.  The Morita evaluation isomorphisms are
\[
 U\otimes_A V\simeq A',\qquad
 V\otimes_{A'}U\simeq A.
\]
Tensor them with $K$.  The canonical base-change isomorphism
\[
 (K\otimes_k U)\otimes_{K\otimes_k A}(K\otimes_k V)
   \simeq K\otimes_k(U\otimes_A V)
\]
and its counterpart with $U,V$ interchanged give the two Morita
evaluation isomorphisms for the extended bimodules.  Their compatibility
with the bimodule actions and with one another is preserved by tensoring.
Thus the extended bimodules give inverse $K$-linear tensor equivalences.
\end{proof}

\begin{proof}[Proof of Theorem~\ref{thm:donovan}]
First take $k=\overline{\F}_p$.  Proposition~\ref{prop:sylow-comparison}
supplies the hypothesis of Theorem~\ref{thm:extension-finiteness}.
Together with Theorem~\ref{thm:quasisimple-cartan}, it proves finiteness
over $k$.  The integral comparisons in this section serve the finite-list
argument for the base blocks; the final scalar normalization, and hence
the assertion for arbitrary crossed extensions, is over $k$.

Now fix an algebraically closed field $K$ of characteristic $p$ and
choose an embedding $k\hookrightarrow K$.  For the fixed bound $M$,
choose blocks $A_1,\ldots,A_t$ representing the finitely many
$k$-linear Morita classes of blocks with defect-group order at most $M$.
Every block $B$ of a finite group algebra $KG$ is, by
Lemma~\ref{lem:coefficient-extension}, of the form
$B=K\otimes_k B_0$ for a unique block $B_0$ of $kG$, and $B_0$ has the
same defect groups as $B$.  If their order is at most $M$, then $B_0$ is
$k$-linearly Morita equivalent to some $A_i$.  The same lemma extends
this equivalence to a $K$-linear Morita equivalence between $B$ and
$K\otimes_k A_i$.  For this fixed field $K$, all required blocks are
therefore represented by the finite list
$K\otimes_k A_1,\ldots,K\otimes_k A_t$.
\end{proof}
